\chapter{2010年湖北卷（文）}

\section{单选题}

\begin{problem}
设集合 \(M = \{1, 2, 4, 8\}\)，\(N = \{x \mid x\text{是 2 的倍数}\}\)，则 \(M \cap N =\)（\quad）

\choices
    {\(\{2, 4\}\)}
    {\(\{1, 2, 4\}\)}
    {\(\{2,4,8\}\)}
    {\(\{1, 2, 4, 8\}\)}
\end{problem}

\begin{answer}
C
\end{answer}

\begin{solution}
集合 \(N\) 中的元素是 \(2\) 的倍数，所以 \(M\) 中同时属于 \(N\) 的元素为 \(2\)，\(4\)，\(8\)，即 \(M\cap N=\{2,4,8\}\). 因此选 C.
\end{solution}

\begin{problem}
函数 \( f(x) = \sqrt{3} \sin \left( \frac{x}{2} - \frac{\pi}{4} \right) \)，\( x \in \R \) 的最小正周期为（\quad）

\choices
    {\(\frac{\pi}{2}\)}
    {\(\pi\)}
    {\(2\pi\)}
    {\(4\pi\)}
\end{problem}

\begin{answer}
D
\end{answer}

\begin{solution}
正弦函数的自变量为 \(\frac{x}{2}-\frac{\pi}{4}\)，其中 \(x\) 的系数为 \(\frac{1}{2}\)，所以最小正周期为 \(T=\frac{2\pi}{1/2}=4\pi\). 因此选 D.
\end{solution}

\begin{problem}
已知函数 \( f(x) = \begin{cases}
\log_{3} x, & x > 0,\\
2^x, & x \leqslant 0,
\end{cases} \)，则 \( f\left(f\left(\frac{1}{9}\right)\right) = \)（\quad）

\choices
    {\(4\)}
    {\(\frac{1}{4}\)}
    {\(-4\)}
    {\(-\frac{1}{4}\)}
\end{problem}

\begin{answer}
B
\end{answer}

\begin{solution}
因为 \(\frac{1}{9}>0\)，所以先使用第一段函数，得 \(f\left(\frac{1}{9}\right)=\log_3\frac{1}{9}=-2\). 又因为 \(-2\leqslant0\)，所以 \(f\left(f\left(\frac{1}{9}\right)\right)=f(-2)=2^{-2}=\frac{1}{4}\). 因此选 B.
\end{solution}

\begin{problem}
用 \(a\)，\(b\)，\(c\) 表示三条不同的直线，\(\gamma\) 表示平面，给出下列命题：

\begin{circlelist}
\item 若 \(a \parallel b\)，\(b \parallel c\)，则 \(a \parallel c\).
\item 若 \(a \perp b\)，\(b \perp c\)，则 \(a \perp c\).
\item 若 \(a \parallel \gamma\)，\(b \parallel \gamma\)，则 \(a \parallel b\).
\item 若 \(a \perp \gamma\)，\(b \perp \gamma\)，则 \(a \parallel b\).
\end{circlelist}

正确的有（\quad）

\choices
    {\circled{1}\circled{2}}
    {\circled{2}\circled{3}}
    {\circled{1}\circled{4}}
    {\circled{3}\circled{4}}
\end{problem}

\begin{answer}
C
\end{answer}

\begin{solution}
命题\(\circled{1}\)成立，因为平行于同一直线的两条直线互相平行. 命题\(\circled{2}\)不成立，例如取 \(b\) 为 \(x\) 轴，取过原点且方向向量分别为 \((0,1,0)\)、\((0,1,1)\) 的两条直线为 \(a\)、\(c\)，则 \(a\perp b\)、\(b\perp c\)，但 \(a\) 与 \(c\) 不垂直. 命题\(\circled{3}\)不成立，例如在平面 \(\gamma:z=0\) 外取两条方向分别平行于 \(x\) 轴、\(y\) 轴的直线，它们都平行于 \(\gamma\)，但彼此不平行. 命题\(\circled{4}\)成立，因为垂直于同一平面的两条直线互相平行. 所以正确的命题是\(\circled{1}\)\(\circled{4}\)，选 C.
\end{solution}

\begin{problem}
函数 \( y = \frac{1}{\sqrt{\log_{0.5}(4x - 3)}} \) 的定义域为（\quad）

\choices
    {\(\left(\frac{3}{4}, 1\right)\)}
    {\(\left(\frac{3}{4}, +\infty\right)\)}
    {\((1, +\infty)\)}
    {\(\left(\frac{3}{4}, 1\right) \cup (1, +\infty)\)}
\end{problem}

\begin{answer}
A
\end{answer}

\begin{solution}
要使对数有意义，需 \(4x-3>0\). 又因为根式在分母中，需 \(\log_{0.5}(4x-3)>0\). 由于 \(0<0.5<1\)，对数函数单调递减，因此 \(0<4x-3<1\)，解得 \(\frac{3}{4}<x<1\). 所以定义域为 \(\left(\frac{3}{4},1\right)\)，选 A.
\end{solution}

\begin{problem}
现有 6 名同学去听同时进行的 5 个课外知识讲座，每名同学可自由选择其中的一个讲座，不同选法的种数是（\quad）

\choices
    {\(5^{6}\)}
    {\(6^{5}\)}
    {\(\frac{5\times 6\times 5\times 4\times 3\times 2}{2}\)}
    {\(6 \times 5 \times 4 \times 3 \times 2\)}
\end{problem}

\begin{answer}
A
\end{answer}

\begin{solution}
每名同学都有 \(5\) 种选择，且 \(6\) 名同学的选择相互独立，所以不同选法的种数为 \(5^6\). 因此选 A.
\end{solution}

\begin{problem}
已知等比数列 \( \{a_{n}\} \) 中，各项都是正数，且 \(a_{1}\)，\(\frac{1}{2}a_{3}\)，\(2a_{2}\) 成等差数列，则 \( \frac{a_{9} + a_{10}}{a_{7} + a_{8}} = \)（\quad）

\choices
    {\(1 + \sqrt{2}\)}
    {\(1 - \sqrt{2}\)}
    {\(3 + 2\sqrt{2}\)}
    {\(3 - 2\sqrt{2}\)}
\end{problem}

\begin{answer}
C
\end{answer}

\begin{solution}
设等比数列的公比为 \(q\). 由三项成等差数列，得 \(2\cdot\frac{1}{2}a_3=a_1+2a_2\)，即 \(a_3=a_1+2a_2\). 因为 \(a_2=a_1q\)，\(a_3=a_1q^2\)，且 \(a_1>0\)，所以 \(q^2=1+2q\)，即 \(q^2-2q-1=0\). 解得 \(q=1\pm\sqrt{2}\)，又各项为正数，故 \(q=1+\sqrt{2}\). 因此
\[
\frac{a_9+a_{10}}{a_7+a_8}
=\frac{a_7q^2+a_7q^3}{a_7+a_7q}
=q^2
=(1+\sqrt{2})^2
=3+2\sqrt{2}
\]
所以选 C.
\end{solution}

\begin{problem}
已知 \(\triangle ABC\) 和点 \(M\) 满足 \(\overrightarrow{MA} + \overrightarrow{MB} + \overrightarrow{MC} = 0\). 若存在实数 \(m\) 使得 \(\overrightarrow{AB} + \overrightarrow{AC} = m\overrightarrow{AM}\) 成立，则 \(m =\)（\quad）

\choices
    {\(2\)}
    {\(3\)}
    {\(4\)}
    {\(5\)}
\end{problem}

\begin{answer}
B
\end{answer}

\begin{solution}
取任意原点，记点 \(A\)，\(B\)，\(C\)，\(M\) 的位置向量分别为 \(\bs{a}\)，\(\bs{b}\)，\(\bs{c}\)，\(\bs{m}\). 题设给出
\[
(\bs{a}-\bs{m})+(\bs{b}-\bs{m})+(\bs{c}-\bs{m})=0
\]
所以 \(\bs{a}+\bs{b}+\bs{c}=3\bs{m}\). 另一方面，
\[
\overrightarrow{AB}+\overrightarrow{AC}
=(\bs{b}-\bs{a})+(\bs{c}-\bs{a})
=\bs{b}+\bs{c}-2\bs{a}
=3\bs{m}-3\bs{a}
=3\overrightarrow{AM}
\]
故 \(m=3\)，选 B.
\end{solution}

\begin{problem}
若直线 \( y = x + b \) 与曲线 \( y = 3 - \sqrt{4x - x^2} \) 有公共点，则 \( b \) 的取值范围是（\quad）

\choices
    {\([1 - 2\sqrt{2}, 1 + 2\sqrt{2}]\)}
    {\([1 - \sqrt{2}, 3]\)}
    {\([-1, 1 + 2\sqrt{2}]\)}
    {\([1 - 2\sqrt{2}, 3]\)}
\end{problem}

\begin{answer}
D
\end{answer}

\begin{solution}
曲线的定义域为 \(0\leqslant x\leqslant4\). 公共点满足
\[
b=3-x-\sqrt{4x-x^2}
\]
记 \(\varphi(x)=3-x-\sqrt{4x-x^2}\). 在 \(0<x<4\) 时，
\[
\varphi'(x)=-1-\frac{2-x}{\sqrt{4x-x^2}}
\]
令 \(\varphi'(x)=0\). 当 \(0<x\leqslant2\) 时，\(2-x\geqslant0\)，故 \(\varphi'(x)<0\)；当 \(2<x<4\) 时，方程 \(\varphi'(x)=0\) 等价于 \(\sqrt{4x-x^2}=x-2\)，平方得 \(4x-x^2=(x-2)^2\)，即 \(x^2-4x+2=0\)，所以 \(x=2+\sqrt{2}\). 结合端点值 \(\varphi(0)=3\)，\(\varphi(4)=-1\)，可得
\(\varphi(2+\sqrt{2})=1-2\sqrt{2}\)，\(\min\varphi=1-2\sqrt{2}\)，\(\max\varphi=3\).
所以 \(b\in[1-2\sqrt{2},3]\)，选 D.
\end{solution}

\begin{problem}
记实数 \(x_{1}\)，\(x_{2}\)，\(\cdots\)，\(x_{n}\) 中的最大数为 \(\max \{x_{1}, x_{2}, \cdots, x_{n}\}\)，最小数为 \(\min \{x_{1}, x_{2}, \cdots, x_{n}\}\). 已知 \(\triangle ABC\) 的三边边长为 \(a\)，\(b\)，\(c\)（\(a \leqslant b \leqslant c\)），定义它的倾斜度为 \(t = \max \left\{\frac{a}{b}, \frac{b}{c}, \frac{c}{a}\right\} \cdot \min \left\{\frac{a}{b}, \frac{b}{c}, \frac{c}{a}\right\}\)，则 “\(t = 1\)” 是 “\(\triangle ABC\) 为等边三角形” 的（\quad）

\choices
    {必要而不充分的条件}
    {充分而不必要的条件}
    {充要条件}
    {既不充分也不必要的条件}
\end{problem}

\begin{answer}
A
\end{answer}

\begin{solution}
由 \(a\leqslant b\leqslant c\)，有 \(\frac{a}{b}\leqslant1\)，\(\frac{b}{c}\leqslant1\)，\(\frac{c}{a}\geqslant1\)，因此最大值为 \(\frac{c}{a}\). 若 \(a=b=c\)，则 \(t=1\)，所以 \(t=1\) 是等边三角形的必要条件.

另一方面，若 \(\frac{a}{b}\leqslant\frac{b}{c}\)，则 \(t=\frac{c}{a}\cdot\frac{a}{b}=\frac{c}{b}\)，此时 \(t=1\) 可推出 \(b=c\)；若 \(\frac{b}{c}\leqslant\frac{a}{b}\)，则 \(t=\frac{c}{a}\cdot\frac{b}{c}=\frac{b}{a}\)，此时 \(t=1\) 可推出 \(a=b\). 因此 \(a=b<c\) 或 \(a<b=c\) 时也可能有 \(t=1\). 例如边长为 \(1\)，\(1\)，\(\frac{3}{2}\) 的三角形满足 \(t=1\)，但不是等边三角形. 所以该条件必要而不充分，选 A.
\end{solution}

\section{填空题}

\begin{problem}
在 \((1 - x^{2})^{10}\) 的展开中，\(x^{4}\) 的系数为\fillinblank{}.
\end{problem}

\begin{answer}
\(45\)
\end{answer}

\begin{solution}
展开式的一般项为 \(\mathrm C_{10}^{k}(-x^2)^k\). 要得到 \(x^4\)，必须有 \(2k=4\)，所以 \(k=2\). 因此 \(x^4\) 的系数为 \(\mathrm C_{10}^{2}(-1)^2=45\).
\end{solution}

\begin{problem}
已知 \(z = 2x - y\)，式中变量 \(x\)，\(y\) 满足约束条件 \(\begin{cases}
y \leqslant x,\\
x + y \geqslant 1,\\
x \leqslant 2
\end{cases}\)，则 \(z\) 的最大值是\fillinblank{}.
\end{problem}

\begin{answer}
\(5\)
\end{answer}

\begin{solution}
由 \(x+y\geqslant1\)，得 \(y\geqslant1-x\). 因为 \(z=2x-y\) 随 \(y\) 的减小而增大，所以在固定 \(x\) 时取 \(y=1-x\). 此时还需满足 \(y\leqslant x\)，即 \(1-x\leqslant x\)，从而 \(x\geqslant\frac12\). 结合 \(x\leqslant2\)，有
\[
z\leqslant2x-(1-x)=3x-1\leqslant5
\]
当 \(x=2\)，\(y=-1\) 时满足全部约束，且 \(z=5\)，所以最大值为 \(5\).
\end{solution}

\begin{problem}
一个病人服用某种新药后被治愈的概率为 \(0.9\). 则服用这种新药的 \(4\) 个病人中至少 \(3\) 人被治愈的概率为\fillinblank{}.
\end{problem}

\begin{answer}
\(0.9477\)
\end{answer}

\begin{solution}
“至少 \(3\) 人被治愈”包括恰有 \(3\) 人被治愈和 \(4\) 人都被治愈两种互斥情况，所以所求概率为
\[
\mathrm C_{4}^{3}(0.9)^3(1-0.9)+(0.9)^4
=4\times0.9^3\times0.1+0.9^4
=0.2916+0.6561
=0.9477
\]
\end{solution}

\begin{problem}
圆柱形容器内部盛有高度为 \(8\mathrm{~cm}\) 的水，若放入三个相同的球（球的半径与圆柱的底面半径相同）后，水恰好淹没最上面的球（如图所示），则球的半径是\fillinblank{}\(\mathrm{cm}\).
\FigureLayoutDeclare{img/2010/hubei_liberal/q14_fig1.png}{3.0cm}{0.8bp 0.8bp 0.8bp 0.8bp}
\bitmapfigure[max width=0.40\linewidth]{img/2010/hubei_liberal/q14_fig1.png}
\end{problem}

\begin{answer}
\(4\)
\end{answer}

\begin{solution}
设球的半径以及圆柱的底面半径均为 \(r\). 三个球放入圆柱后，最上面的球恰好被水淹没，所以水面高度为 \(6r\). 放入球前水的体积为 \(8\pi r^2\)，放入球后水的体积等于高为 \(6r\) 的圆柱体积减去三个球的体积，即
\[
8\pi r^2=6\pi r^3-3\cdot\frac{4}{3}\pi r^3=2\pi r^3
\]
因为 \(r>0\)，解得 \(r=4\). 所以球的半径为 \(4\mathrm{~cm}\).
\end{solution}

\begin{problem}
已知椭圆 \(C: \frac{x^2}{2} + y^2 = 1\) 的两焦点为 \(F_1\)，\(F_2\)，点 \(P(x_0, y_0)\) 满足 \(0 < \frac{x_0^2}{2} + y_0^2 < 1\)，则 \(|PF_1| + |PF_2|\) 的取值范围为\fillinblank{}，直线 \(\frac{x_0 x}{2} + y_0 y = 1\) 与椭圆 \(C\) 的公共点个数为\fillinblank{}.
\end{problem}

\begin{answer}
\([2,2\sqrt{2})\)，\(0\)
\end{answer}

\begin{solution}
椭圆 \(C\) 的长半轴为 \(\sqrt{2}\)，焦半距为 \(1\)，所以椭圆上任一点到两焦点距离之和为 \(2\sqrt{2}\)，椭圆内部任一点到两焦点距离之和小于 \(2\sqrt{2}\). 又由三角不等式，
\[
|PF_1|+|PF_2|\geqslant|F_1F_2|=2
\]
当 \(P\) 在焦点连线段 \(F_1F_2\) 上时等号成立，且这类点中除原点外均满足题设条件，因此取值范围为 \([2,2\sqrt{2})\).

令 \(u=\left(\frac{x_0}{\sqrt{2}},y_0\right)\)，\(v=\left(\frac{x}{\sqrt{2}},y\right)\). 椭圆上的点满足 \(|v|=1\)，而题设给出 \(|u|=\sqrt{\frac{x_0^2}{2}+y_0^2}<1\). 直线方程要求 \(u\cdot v=1\)，但由柯西不等式 \(u\cdot v\leqslant|u||v|=|u|<1\)，矛盾. 因此公共点个数为 \(0\).
\end{solution}

\section{解答题}

\begin{problem}
已知函数 \( f(x) = \frac{\cos^2 x - \sin^2 x}{2} \)，\( g(x) = \frac{1}{2} \sin 2x - \frac{1}{4} \).
\begin{enumerate}
\item 函数 \( f(x) \) 的图象可由函数 \( g(x) \) 的图象经过怎样变化得出？
\item 求函数 \( h(x)=f(x)-g(x) \) 的最小值，并求使 \( h(x) \) 取得最小值的 \(x\) 的集合.
\end{enumerate}
\end{problem}

\begin{solution}
\begin{enumerate}
\item 由 \(\cos^2x-\sin^2x=\cos2x\)，得 \(f(x)=\frac12\cos2x\). 又
\[
g\left(x+\frac{\pi}{4}\right)+\frac14
=\frac12\sin\left(2x+\frac{\pi}{2}\right)
=\frac12\cos2x
=f(x)
\]
所以 \(f(x)\) 的图象可由 \(g(x)\) 的图象向左平移 \(\frac{\pi}{4}\) 个单位，再向上平移 \(\frac14\) 个单位得到.
\item
\[
h(x)=\frac12\cos2x-\frac12\sin2x+\frac14
=\frac{\sqrt{2}}{2}\cos\left(2x+\frac{\pi}{4}\right)+\frac14
\]
因此 \(h(x)\) 的最小值为 \(\frac14-\frac{\sqrt{2}}2\). 取得最小值时
\[
\cos\left(2x+\frac{\pi}{4}\right)=-1
\]
即 \(2x+\frac{\pi}{4}=\pi+2k\pi\)，\(k\in\Z\). 所以 \(x=\frac{3\pi}{8}+k\pi\)，\(k\in\Z\)，所求集合为 \(\left\{\frac{3\pi}{8}+k\pi\mid k\in\Z\right\}\).
\end{enumerate}
\end{solution}

\begin{problem}
为了了解一个小水库中养殖的鱼的有关情况，从这个水库中多个不同位置捕捞出 \(100\) 条鱼，称得每条鱼的质量（单位：千克），并将所得数据分组，画出频率分布直方图（如图所示）.

\begin{center}
\begin{tabular}{c|c}
\hline
分组 & 频率 \\
\hline
\([1.00,1.05)\) & \\
\([1.05,1.10)\) & \\
\([1.10,1.15)\) & \\
\([1.15,1.20)\) & \\
\([1.20,1.25)\) & \\
\([1.25,1.30)\) & \\
\hline
\end{tabular}
\end{center}

\begin{enumerate}
\item 在下面表格中填写相应的频率；
\item 估计数据落在 \([1.15, 1.30)\) 中的概率为多少；
\item 将上面捕捞的 \(100\) 条鱼分别作一记号后再放回水库，几天后再从水库的多次不同位置捕出 \(120\) 条鱼，其中带有记号的鱼有 \(6\) 条，请根据这一情况来估计该水库中鱼的总条数.
\end{enumerate}

\texfigure[max width=0.72\linewidth]{img_repaint/2010/hubei_liberal/q17_fig1.tex}
\end{problem}

\begin{solution}
\begin{enumerate}
\item 直方图的纵坐标表示频率除以组距，各组组距均为 \(0.05\). 由图可知六组的频率密度依次为 \(1\)，\(4\)，\(5.6\)，\(6\)，\(3\)，\(0.4\)，所以相应频率依次为
\(1\times0.05=0.05\)，\(4\times0.05=0.20\)，\(5.6\times0.05=0.28\)，\(6\times0.05=0.30\)，\(3\times0.05=0.15\)，\(0.4\times0.05=0.02\).
这些频率之和为 \(1\)，与频率分布相符.
\item 数据落在 \([1.15,1.30)\) 内的频率为 \(0.30+0.15+0.02=0.47\)，所以所估计的概率为 \(0.47\).
\item 设水库中鱼的总数为 \(N\). 用样本中带记号的鱼所占比例估计水库中带记号鱼所占比例，有 \(\frac{100}{N}\approx\frac{6}{120}\)，所以 \(N\approx\frac{100\times120}{6}=2000\). 因此估计该水库中鱼的总数约为 \(2000\) 条.
\end{enumerate}
\end{solution}

\begin{problem}
如图，在四面体 \(ABOC\) 中，\(OC \perp OA\)，\(OC \perp OB\)，\(\angle AOB = 120^\circ\)，且 \(OA = OB = OC = 1\).
\begin{enumerate}
\item 设 \(P\) 为 \(AC\) 的中点，\(Q\) 在 \(AB\) 上且 \(AB = 3AQ\)，证明： \(PQ \perp OA\)；
\item 求二面角 \(O - AC - B\) 的平面角的余弦值.
\end{enumerate}
\bitmapfigure[max width=0.44\linewidth]{img/2010/hubei_liberal/q18_fig1.png}
\end{problem}

\begin{solution}
\begin{enumerate}
\item 取 \(O\) 为原点，令 \(OA\) 为 \(x\) 轴正向，\(OC\) 为 \(z\) 轴正向. 由题设可取 \(A=(1,0,0)\)，\(C=(0,0,1)\)，\(B=\left(-\frac12,\frac{\sqrt{3}}2,0\right)\)，因为 \(\lvert OA\rvert=\lvert OC\rvert=1\)，\(\lvert OB\rvert=\sqrt{\frac14+\frac34}=1\)，\(\overrightarrow{OA}\cdot\overrightarrow{OC}=\overrightarrow{OB}\cdot\overrightarrow{OC}=0\)，且 \(\overrightarrow{OA}\cdot\overrightarrow{OB}=-\frac12=\lvert OA\rvert\lvert OB\rvert\cos120^\circ\). 因为 \(P\) 是 \(AC\) 的中点，所以 \(P=\left(\frac12,0,\frac12\right)\). 又 \(AQ=\frac13AB\)，故
\[
Q=A+\frac13(B-A)=\left(\frac12,\frac{\sqrt{3}}6,0\right)
\]
于是 \(\overrightarrow{PQ}=\left(0,\frac{\sqrt{3}}6,-\frac12\right)\)，而 \(\overrightarrow{OA}=(1,0,0)\)，所以 \(\overrightarrow{PQ}\cdot\overrightarrow{OA}=0\). 因此 \(PQ\perp OA\).
\item 平面 \(OAC\) 的一个法向量为 \(\bs{n}_1=(0,1,0)\). 取
\(\overrightarrow{AB}=\left(-\frac32,\frac{\sqrt{3}}2,0\right)\)，\(\overrightarrow{AC}=(-1,0,1)\)
则平面 \(ABC\) 的一个法向量为
\[
\bs{n}_2=\overrightarrow{AB}\times\overrightarrow{AC}
=\left(\frac{\sqrt{3}}2,\frac32,\frac{\sqrt{3}}2\right)
\]
因此二面角 \(O-AC-B\) 的平面角 \(\theta\) 满足
\[
\cos\theta=\frac{\bs{n}_1\cdot\bs{n}_2}{|\bs{n}_1||\bs{n}_2|}
=\frac{3/2}{\sqrt{15}/2}
=\frac{\sqrt{15}}5
\]
\end{enumerate}
\end{solution}

\begin{problem}
已知一条曲线 \(C\) 在 \(y\) 轴的右边，\(C\) 上每一点到点 \(F(1,0)\) 的距离减去它到 \(y\) 轴距离的差都是 \(1\).
\begin{enumerate}
\item 求曲线 \(C\) 的方程；
\item 是否存在正数 \(m\)，对于过点 \(M(m,0)\) 且与曲线 \(C\) 有两个交点 \(A\)，\(B\) 的任一直线，都有 \(\overrightarrow{FA} \cdot \overrightarrow{FB} < 0\)？若存在，求出 \(m\) 的取值范围；若不存在，请说明理由.
\end{enumerate}
\end{problem}

\begin{solution}
\begin{enumerate}
\item 设 \(C\) 上任一点为 \(P(x,y)\). 因为 \(P\) 在 \(y\) 轴右边，所以 \(x>0\)，且 \(P\) 到 \(F(1,0)\) 和 \(y\) 轴的距离分别为 \(\sqrt{(x-1)^2+y^2}\) 和 \(x\). 由题意，
\[
\sqrt{(x-1)^2+y^2}-x=1
\]
于是 \(\sqrt{(x-1)^2+y^2}=x+1\)，两边均非负，可以平方得
\[
(x-1)^2+y^2=(x+1)^2
\]
即 \(y^2=4x\). 所以曲线 \(C\) 的方程为 \(y^2=4x\)，并取 \(x>0\) 的部分.
\item 设直线与抛物线交于参数点 \(A=(u^2,2u)\)，\(B=(v^2,2v)\)，其中 \(u\neq v\). 当直线过 \(M(m,0)\) 时，由弦的方程或直接代入三点共线条件可得 \(uv=-m\). 这对竖直弦也成立，因为此时 \(u+v=0\)，且 \(m=u^2=-uv\).

由 \(F=(1,0)\)，有
\(\overrightarrow{FA}=(u^2-1,2u)\)，\(\overrightarrow{FB}=(v^2-1,2v)\).
记 \(s=u+v\)，则 \(uv=-m\)，从而
\[
\overrightarrow{FA}\cdot\overrightarrow{FB}=(u^2-1)(v^2-1)+4uv=u^2v^2-u^2-v^2+1+4uv=m^2-(s^2+2m)+1-4m=m^2-6m+1-s^2
\]
对过 \(M\) 的任意割线，\(s\) 可以取任意实数，特别地，竖直直线给出 \(s=0\). 事实上，任取 \(s\in\R\)，方程 \(\lambda^2-s\lambda-m=0\) 的判别式为 \(s^2+4m>0\)，其两个不同实根 \(u\)、\(v\) 满足 \(u+v=s\)、\(uv=-m<0\)，从而对应曲线 \(C\) 上 \(x>0\) 的两个交点. 因而要使上述内积对任意割线都严格小于 \(0\)，必须且只需 \(m^2-6m+1<0\). 解得
\[
3-2\sqrt{2}<m<3+2\sqrt{2}
\]
这个区间中的数均为正，所以满足条件的 \(m\) 的取值范围为 \((3-2\sqrt{2},3+2\sqrt{2})\).
\end{enumerate}
\end{solution}

\begin{problem}
设函数 \( f(x) = \frac{1}{3} x^3 - \frac{a}{2} x^2 + bx + c \)，其中 \( a > 0 \)，曲线 \( y = f(x) \) 在点 \( P(0, f(0)) \) 处的切线方程为 \( y = 1 \).
\begin{enumerate}
\item 确定 \(b\)，\(c\) 的值；
\item 设曲线 \( y = f(x) \) 在点 \( (x_1, f(x_1)) \) 及 \( (x_2, f(x_2)) \) 处的切线都过点 \((0,2)\). 证明：当 \( x_1 \neq x_2 \) 时，\( f'(x_1) \neq f'(x_2)\).
\item 若过点 \((0,2)\) 可作曲线 \( y = f(x) \) 的三条不同切线，求 \(a\) 的取值范围.
\end{enumerate}
\end{problem}

\begin{solution}
\begin{enumerate}
\item \(f(0)=c\)，且 \(f'(x)=x^2-ax+b\). 切线 \(y=1\) 经过 \(P(0,f(0))\)，所以 \(c=1\)；该切线斜率为 \(0\)，所以 \(f'(0)=b=0\).
\item 由第（1）问，\(f(x)=\frac13x^3-\frac a2x^2+1\)，\(f'(x)=x^2-ax\). 曲线在 \(x=t\) 处的切线过 \((0,2)\) 的充要条件为
\[
f(t)-tf'(t)=2
\]
即
\[
1+\frac13t^3-\frac a2t^2-t(t^2-at)=2
\]
化简得 \(4t^3-3at^2+6=0\). 因而 \(x_1\)，\(x_2\) 都是方程 \(4x^3-3ax^2+6=0\) 的根.

若 \(f'(x_1)=f'(x_2)\)，则
\[
(x_1-x_2)(x_1+x_2-a)=0
\]
由于 \(x_1\neq x_2\)，可得 \(x_1+x_2=a\). 又因为 \(x_1\)，\(x_2\) 都是上述方程的根，作差并除以 \(x_1-x_2\)，得到
\[
4(x_1^2+x_1x_2+x_2^2)-3a(x_1+x_2)=0
\]
代入 \(x_1+x_2=a\)，可得 \(x_1x_2=\frac{a^2}{4}\). 于是 \(x_1\)，\(x_2\) 是方程 \(X^2-aX+\frac{a^2}{4}=0\) 的两个根，但该方程为 \(\left(X-\frac a2\right)^2=0\)，从而 \(x_1=x_2\)，矛盾. 所以 \(f'(x_1)\neq f'(x_2)\).
\item 令 \(\phi(x)=4x^3-3ax^2+6\). 则
\[
\phi'(x)=6x(2x-a)
\]
因为 \(a>0\)，所以 \(\phi\) 在 \((-\infty,0)\) 上递增，在 \((0,a/2)\) 上递减，在 \((a/2,+\infty)\) 上递增. 且
\(\phi(0)=6\)，\(\phi\left(\frac a2\right)=6-\frac{a^3}{4}\).
又 \(\lim_{x\to-\infty}\phi(x)=-\infty\)，\(\lim_{x\to+\infty}\phi(x)=+\infty\)，所以方程 \(\phi(x)=0\) 有三个不同实根，当且仅当局部极小值小于 \(0\)，即 \(6-\frac{a^3}{4}<0\). 因此 \(a^3>24\)，即 \(a>\sqrt[3]{24}\). 由第（2）问，三个不同根对应三条不同切线，所以 \(a\) 的取值范围为 \((\sqrt[3]{24},+\infty)\).
\end{enumerate}
\end{solution}

\begin{problem}
已知某地今年年初拥有居民住房的总面积为 \(a\)（单位：\(\mathrm{m}^2\)），其中有部分旧住房需要拆除. 当地有关部门决定每年以当年年初住房面积的 \(10\%\) 建设新住房，同时也拆除面积为 \(b\)（单位：\(\mathrm{m}^2\)）的旧住房.
\begin{enumerate}
\item 分别写出第一年末和第二年末的实际住房面积的表达式；
\item 如果第五年末该地的住房面积正好比今年年初的住房面积增加了 \(30\%\)，则每年拆除的旧住房面积 \(b\) 是多少？（计算时取 \(1.1^{5}=1.6\)）
\end{enumerate}
\end{problem}

\begin{solution}
\begin{enumerate}
\item 设第 \(n\) 年年初的住房面积为 \(S_{n-1}\)，则每年建设后并拆除旧住房，年末面积满足 \(S_n=1.1S_{n-1}-b\)，且 \(S_0=a\). 因此第一年末
\[
S_1=1.1a-b
\]
第二年末
\[
S_2=1.1S_1-b=1.1(1.1a-b)-b=1.21a-2.1b
\]
\item 由递推关系逐次展开，得第五年末
\[
S_5=1.1^5a-b(1+1.1+1.1^2+1.1^3+1.1^4)
\]
按题设 \(S_5=1.3a\)，且 \(1.1^5=1.6\)，所以
\[
1+1.1+1.1^2+1.1^3+1.1^4=\frac{1.1^5-1}{1.1-1}=\frac{1.6-1}{0.1}=6
\]
于是 \(1.6a-6b=1.3a\)，解得 \(b=0.05a\).
\end{enumerate}
\end{solution}
